% ---------------------------------------------------------------------
%  CHƯƠNG 5 — KẾT QUẢ VÀ ĐÁNH GIÁ
% ---------------------------------------------------------------------
\chapter{Kết quả và đánh giá}
\label{ch:ketqua}

\begin{muctieu}
Chương này trình bày kết quả đạt được, các độ đo đánh giá định lượng, phần
bàn luận và những hạn chế còn tồn tại.
\end{muctieu}

\section{Kết quả đạt được}

\begin{chuy}
\goiy{đối chiếu trực tiếp từng mục tiêu đã nêu ở Chương~\ref{ch:modau} để
người đọc thấy mức độ hoàn thành. Mục tiêu nào chưa đạt thì nói thẳng, đừng
lờ đi.}
\end{chuy}

\begin{mau}
Hệ thống đã hoàn thành các mục tiêu đề ra. Mô hình phân loại đạt độ chính
xác vượt ngưỡng 80\% trên tập kiểm thử, quy trình tiền xử lý hoạt động ổn
định và giao diện web cho phép dự đoán theo thời gian thực.
\end{mau}

\section{Độ đo đánh giá}

\begin{chuy}
\goiy{trình bày các độ đo phù hợp với bài toán, ví dụ độ chính xác,
precision, recall, F1 cho bài toán phân lớp, hoặc MAE và RMSE cho bài toán
dự báo. Ghi rõ số liệu đo trên tập nào và với bao nhiêu mẫu.}
\end{chuy}

Bảng~\ref{tab:ketqua} so sánh hai mô hình trên cùng tập kiểm thử gồm 1.000 mẫu.

\begin{table}[H]
\centering
\caption{So sánh hiệu năng các mô hình trên tập kiểm thử.}
\label{tab:ketqua}
\begin{tabular*}{\textwidth}{@{\extracolsep{\fill}}>{\bfseries\color{daunavy}}l cccc}
\toprule
\rowcolor{gray100}\bfseries\color{daunavy}Mô hình & \bfseries\color{daumaroon}Accuracy & \bfseries\color{daumaroon}Precision & \bfseries\color{daumaroon}Recall & \bfseries\color{daumaroon}F1 \\
\midrule
Hồi quy logistic  & 0,84 & 0,83 & 0,82 & 0,83 \\
Máy vec-tơ hỗ trợ & 0,86 & 0,85 & 0,84 & 0,85 \\
\bottomrule
\end{tabular*}
\end{table}

\section{Bàn luận}

\begin{chuy}
\goiy{giải thích vì sao kết quả lại như vậy, chứ không chỉ đọc lại bảng số.
Mô hình sai ở những ca nào và vì sao. Đây là phần phân biệt báo cáo khá với
báo cáo tốt.}
\end{chuy}

\begin{mau}
Máy vec-tơ hỗ trợ nhỉnh hơn hồi quy logistic khoảng hai điểm phần trăm ở mọi
độ đo. Khoảng cách này không lớn và phải trả giá bằng thời gian huấn luyện
lâu hơn. Phân tích các ca dự đoán sai cho thấy phần lớn lỗi tập trung ở nhãn
trung tính, nơi ranh giới ngữ nghĩa với hai nhãn còn lại là mờ nhất.
\end{mau}

\section{Hạn chế}

\begin{chuy}
\goiy{thẳng thắn nêu những điểm hệ thống chưa làm tốt: dữ liệu thiếu, lỗi
với câu phức tạp, chưa tối ưu tốc độ. Phần này thể hiện tư duy phản biện và
luôn được hội đồng ghi nhận.}
\end{chuy}

\begin{luuy}
Mô hình còn nhầm lẫn với các câu mang sắc thái châm biếm và với bình luận
chứa nhiều từ viết tắt hoặc tiếng lóng. Tập dữ liệu chưa cân bằng giữa các
nhãn cũng làm giảm độ chính xác ở nhãn trung tính.
\end{luuy}
